\section{公式}

\subsection{点电荷系的静电能}
\begin{align}
    \phi =\frac{\rho(r_{source})\rho(r_{acceptor})}{4\pi \epsilon_0(r_{acceptor}-r_{source})^{2}}
\end{align}
\begin{equation}
    \begin{split}
    &\quad \phi(\rho(r_{1}),\rho(r_{2}),\dots,\rho(r_{n}))\\
    &=\frac{\rho(r_{1})\rho(r_{2})}{(r_{2}-r_{1})^{2}}\\
    &\quad +(\frac{\rho(r_{1})\rho(r_{3})}{(r_{3}-r_{1})^{2}}+\frac{\rho(r_{2})\rho(r_{3})}{(r_{3}-r_{2})^{2}})\\
    &\quad +(\frac{\rho(r_{1})\rho(r_{4})}{(r_{4}-r_{1})^{2}}+\frac{\rho(r_{2})\rho(r_{4})}{(r_{4}-r_{2})^{2}}+\frac{\rho(r_{3})\rho(r_{4})}{(r_{4}-r_{3})^{2}})\\
    &\quad +(\frac{\rho(r_{1})\rho(r_{5})}{(r_{5}-r_{1})^{2}}+\frac{\rho(r_{2})\rho(r_{5})}{(r_{5}-r_{2})^{2}}+\frac{\rho(r_{3})\rho(r_{5})}{(r_{5}-r_{3})^{2}}+\frac{\rho(r_{4})\rho(r_{5})}{(r_{5}-r_{4})^{2}})\\
    &\quad +\dots \\
    &=\sum_{j=2}^{n}\sum_{i=1}^{j-1}\frac{\rho(r_{i})\rho(r_{j})}{(r_{j}-r_{i})^{2}}\\
    &=\frac{1}{2}\sum_{j=1}^{n}\sum_{i=1}^{n}\frac{\rho(r_{i})\rho(r_{j})}{(r_{j}-r_{i})^{2}}\\
    \end{split}
\end{equation}
\begin{align}
    \intop_{r=a}^{r=b}{\intop_{r'=a}^{r'=r}{\frac{1}{(r-r')^{2}}dr'dr}}
\end{align}

电流元$Id\vec{l}=\vec{j}dV$，其中，$I=\vec{j}\cdot \vec{A}$, $dV=\vec{A}\cdot d\vec{l}$，
$\vec{A}$为电流穿过的横截面面积向量。$d\vec{l}=dl\vec{n}_{0}$与$\vec{j}=j\vec{n}_{0}$同向，
$Id\vec{l}=(\vec{j}\cdot \vec{A})d\vec{l}=j\vec{n}_{0}\cdot \vec{A}dl\vec{n}_{0}$，
$\vec{j}dV=\vec{j}(\vec{A}\cdot d\vec{l})=j\vec{n}_{0}(\vec{A}\cdot dl\vec{n}_{0})$。

\par
$dV=Sdl\ v.s.\ dr^{3}$

\subsection{构造法：比奥萨伐尔定律的启发 (没对)}
已知：1）$F$的具体形式，2）$F$是某个矢量场$A$的旋度场，$F=\nabla\times A$，求：$A$。

物理上，我们知道恒稳磁场$\nabla\times \vec{B}=\mu_0\vec{j}$，其中$\vec{j}$是电流密度。
而磁感应强度$B$又可以用电流$j$计算出来，由比奥萨伐尔定律
$\vec{B}(\vec{r_{2}}) = \frac{\mu_0}{4\pi} \int \frac{\vec{j}(\vec{r_{1}}) \times (\vec{r_{2}}-\vec{r_{1}})}{|\vec{r_{2}}-\vec{r_{1}}|^3} dV_{1}$

也即，$\vec{F}=\nabla\times \vec{A}$，
$\vec{A}(\vec{r_{2}}) = \frac{\mu_0}{4\pi} \int \frac{\vec{F}(\vec{r_{1}}) \times (\vec{r_{2}}-\vec{r_{1}})}{|\vec{r_{2}}-\vec{r_{1}}|^3} dV_{1}$

\subsection{三维分部积分}
希腊字母代表标量场，大写英文字母代表矢量场，与$\nabla$两两组合：
1）$(\psi\phi)$、$(\psi A)$、$(A\cdot B)$、$(A\times B)$，
2）$\nabla \psi$、$\nabla\cdot A$、$\nabla\times A$，
3）$\nabla\cdot\nabla=\nabla^2$、$\nabla \times \nabla$

1个$\nabla$与（2个场）组合（对应分部积分）：

$\nabla (\psi\phi)$

$\nabla\cdot (\psi A)$、$\nabla\times (\psi A)$

$\nabla (A\cdot B)$

$\nabla \cdot (A\times B)$、$\nabla\times (A\times B)$

1个场与（1个$\nabla$+1个场）组合：

2个$\nabla$与1个场组合：

\subsection{波函数、格林函数法}

$\phi\varphi$

波函数$\psi(\alpha_{i})$只表示一个量子态的态矢量$|\psi\rangle$在$\alpha$表象下投影到的基矢$|\alpha_i\rangle$的\textbf{分量}，
$\psi(\alpha_{i})=\langle\alpha_i|\psi\rangle$，它本身甚至都不是一个态矢量（$|\alpha_i\rangle\langle\alpha_i|\psi\rangle=\psi(\alpha_{i})|\alpha_i\rangle$才是）。

完备性关系$\sum_i |\alpha_i\rangle\langle\alpha_i| = I$可以将量子态$|\psi\rangle$在$\alpha$表象上完整地表示出来，
$|\psi\rangle=\sum_i |\alpha_i\rangle\langle\alpha_i|\psi\rangle=\sum_i \psi(\alpha_i)|\alpha_i\rangle $。

一组构成完备基的可观测量=一个表象，常见的可观测量有：位置$x$、动量$p$、能量（动能、势能）$E$、角动量$L$、自旋$S$、宇称$P$、电荷共轭$C$、粒子数$N$等。

$\hat{O}|\psi\rangle=V|\psi\rangle$，$\hat{O}|\phi\rangle=0$

解：

$|\psi\rangle=|\phi\rangle+\hat{O}^{-1}V|\psi\rangle$

$\langle x|\psi\rangle=\langle x|\phi\rangle+\langle x|\hat{O}^{-1}V|\psi\rangle$

$\langle x|\psi\rangle=\langle x|\phi\rangle+\langle x|\hat{O}^{-1}(\sum |x'\rangle\langle x'|)V|\psi\rangle$

$\langle x|\psi\rangle=\langle x|\phi\rangle+\langle x|\hat{O}^{-1}(\int |x'\rangle\langle x'|dx') V|\psi\rangle=\langle x|\phi\rangle+\int \langle x|\hat{O}^{-1}|x'\rangle\langle x'|V|\psi\rangle dx'$

\begin{equation}
    \psi(x)=\phi(x)+\int O^{-1}(x,x') V(x') \psi(x') dx'
    \label{eq:resolvent}
\end{equation}

其中，
1）算符$\hat{O}^{-1}$在$x$表象下的矩阵元$\langle x|\hat{O}^{-1}|x'\rangle$可以写为函数$O^{-1}(x,x')$，
2）$\langle x'|V|\psi\rangle=\langle x'|V(\int |x''\rangle\langle x''|dx'') |\psi\rangle=\int \langle x'|V|x''\rangle \psi(x'')dx''$，
势能$\langle x'|V|x''\rangle=V(x')\delta(x'-x'')$，
$\langle x'|V|\psi\rangle=\int V(x')\delta(x'-x'') \psi(x'') dx''=V(x')\psi(x')$

再看原方程：

$\langle x|\hat{O}|\psi\rangle=\langle x|V|\psi\rangle$

$\int O(x,x')\psi(x')dx'=V(x)\psi(x)$，代入式\eqref{eq:resolvent}第二项：

\begin{equation}
    \begin{split}
    &\quad \int O^{-1}(x,x') V(x') \psi(x') dx'\\
    =& \int O^{-1}(x,x') \int O(x',x'')\psi(x'')dx'' dx'\\
    =& \int (\int O^{-1}(x,x') O(x',x'')dx')\psi(x'') dx''\\
    =& \int \delta(x-x'')\psi(x'') dx''\\
    =& \psi(x)
    \end{split}
\end{equation}
其中，$\int O^{-1}(x,x') O(x',x'')dx'=\delta(x-x'')$是正逆矩阵相乘$\hat{O}\hat{O}^{-1}=I$的矩阵元展开
